%!TEX root = ../template.tex
%%%%%%%%%%%%%%%%%%%%%%%%%%%%%%%%%%%%%%%%%%%%%%%%%%%%%%%%%%%%%%%%%%%%
%% chapter5.tex
%% NOVA thesis document file
%%
%% Chapter with lots of dummy text
%%%%%%%%%%%%%%%%%%%%%%%%%%%%%%%%%%%%%%%%%%%%%%%%%%%%%%%%%%%%%%%%%%%%

\typeout{NT FILE chapter5.tex}%

\chapter{Conclusions}
\label{cha:conclusions}

This thesis aims to implement a fully functional 5G \gls{SA} network, including both \gls{RAN} and Core elements, built entirely with open-source components and deployed on commodity hardware. The main goals are to deploy the network in a testbed environment, logging the necessary steps of development and configurations to allow for reproducibility, and to evaluate its performance under different scenarios, measuring \glspl{KPI}. 

A review of the current landscape shows that many projects around the implementation of open-source 5G networks have been carried out in recent years, employing several different tools and providing extensive documentation and analysis of results. Building upon this existing body of work, this thesis will focus on deploying a 5G testbed tailored to meet the specific requirements of our environment, while also contributing to the broader research community, where possible, by sharing insights and findings.

Moving forward, the next steps of this thesis will follow the methodology outlined in Chapter~\ref{cha:workplan}, starting with the setup of the 5G \gls{RAN} and Core components, followed by the setup and fine-tuning of the \gls{SDR} hardware. Finally, the performance evaluation will be conducted, gathering \glsxtrlongpl{KPI} values and conluding with an analysis of the results obtained.
% [Talk about aim and goals] \\
% [Talk about current landscape] \\
% [Talk about the next step] 